\section{Nuclear Energy, Evolution, the Unconscious}

When science changed its methodology from the description of phenomena, to the search for occult causes, it then made its three most significant discoveries of modern times. These are \textbf{Nuclear Energy} in Physics, \textbf{Evolution} in Biology, and the \textbf{Unconscious} in Psychology.

What they have in common is that the cause cannot be directly observed, but must be inferred from observation of its effects. The other characteristic they have in common is their ultimate unintelligibility, despite the large amount of experimental data to back them up. This fits in with the program of Positivism: collect the facts and derive laws from them. Metaphysical questions are irrelevant.



\paragraph{Energy}


In physics, matter is condensed energy. Furthermore, at the atomic and subatomic levels, energy persists as probability waves. Its condensation into matter is therefore indeterminate. Nevertheless, the statistical distribution of matter can be determined with a high degree of accuracy. So we can conclude that, from a quantitative perspective, quantum theory is extremely well attested.

The discussion is much different when we look at quantum theory qualitatively, that is, the world picture entailed by the theory. The human mind eventually rebels at positivism and looks for intelligibility, or an ulterior explanation. Unfortunately, there are about a dozen competing extra-scientific interpretations\footnote{\url{https://en.wikipedia.org/wiki/Interpretations_of_quantum_mechanics}} of the quantitative success of the theory. Naturally, the more mystical versions, such as the many-worlds interpretation, are the ones that are popularized.


\paragraph{Evolution}


Because of the fossil record and the similarity of DNA between closely related species, the theory of evolution is well attested in biology. But again, there is a lack of intelligibility. As with quantum mechanics, all explanations are stochastic and statistical. Looking backwards, a biologist can speculate about the species linked up to the appearance of man. But there is no explanation as to why man appeared in the first place.


\paragraph{Unconscious}


The final discovery is that of the unconscious in human motivation. By definition, the unconscious is not observable. Freud attributed human behavior to the unconscious energy of instinctive processes. Adler, following Nietzsche, postulated an unconscious will to power. Finally, Jung noted several higher level unconscious processes, which he called archetypes. There is no explanation for these unconscious forces and the psychological techniques for bringing them into awareness are unreliable.

\paragraph{Esoteric interpretation}


These three scientific discoveries dominate modern and postmodern thought:

\begin{itemize}


\item The world is the result of spontaneous configurations of energy

\item Man is the result of random process of evolution

\item Behaviour is dominated by unconscious motives
\end{itemize}


This model leaves no room for human consciousness, and, a fortiori, for any influences from forces transcendent to the world process. Intelligibility means knowledge of the sufficient reason for things and events, and as long as the only principle of sufficient reason acceptable is physical causation, the postmodern worldview is literally unintelligent. Since no one can live by random and unconscious processes alone, intelligibility is replaced by ideology.


\flrightit{Posted on 2010-10-22 by Cologero}

\begin{center}* * *\end{center}

\begin{footnotesize}
\begin{sffamily}
\texttt{James O'Meara on 2010-10-23 at 10:01 said:}

``Intelligibility means knowledge of the sufficient reason for things and events, and as long as the only principle of sufficient reason acceptable is physical causation, the postmodern worldview is literally unintelligent. ''

This is why Evola disparaged ``science'' in favor of esotericism. As expressed most forcefully in various places in ``Magic'', modern science gives predictability, but real Knowledge involves knowing, or rather, becoming, the remote cause. He gives the example of the scientist pushing a button and sending off a rocket to destroy a distant target. This is the man in the street's paradigm of `knowledge,' but as E. points out, the scientist, qua man, has no real control over any of these forces, and also, is unchanged, as a man, from the most intelligent ape. Thus, although E. doesn't say this, the perennial CP Snow problem of modern men as ``monkeys with machine guns''. ends and means, etc. Evola presents esoteric knowledge as adhering to experience, like sense knowledge, but evoking superior experience as its token of certainty, unlike the merely abstract mathematical models of science. Your ``occult causes'' threw me at the start; I think you have the same model, but by ``description of phenomena'' I assume you mean E.'s superior experience, and by ``occult causes'' you mean his abstract models of science.

Also, an interesting way to show how ideology viz, the ``bright idea'' of college professors has gained its grip on modern man.

\hfill

\texttt{Cologero on 2010-10-23 at 13:28 said:}

Yes, there is some irony in ``occult science'', intended in its literal definition as hidden. It alludes to science as a sort of counter-initiation, as it discovers the hidden nature of matter, life, and mind.

It is the opposite of the metaphysical axiom that the \emph{subtle rules the dense}.

Instead of using knowledge of unconscious forces (e.g., the function of astrology) to liberate, it is used as a means of control. For example, all that has followed from Edward Bernays's use of Uncle Sigmund's theories. Even \textit{children's cartoons}\footnote{\url{http://www.howstuffworks.com/how-dora-the-explorer-works.htm/printable}} are reflect the same attention to manipulation.

Similarly, the ideal of the ``true man'' or even Nietzsche's ``superman'', which are conscious outcomes of the human will, will be replaced by genetic manipulation. If anyone thinks this will somehow improve the human race, he needs to take a close look at how the wolf became a ``dog''.

\hfill

\texttt{Daniel on 2010-10-24 at 02:26 said:}

This is very insightful, especially the second half. Modern science, when it encounters the limits of what it can explain, engages in a lot of elaborate hand-waving, which often is sufficient to distract one from the obvious truth, which is: they don't know the answer; the answer is unknowable using their postivist methods.

As a side note, what do you mean by the ``attention to manipulation'' in \textbf{Dora the Explorer}?


\hfill

\texttt{Cologero on 2010-10-24 at 13:38 said:}

I am far from an expert on pop culture, but I had the opportunity recently to witness young children watching Dora the Explorer. First I noted how different the cartoon is from, say, Bugs Bunny, yet the children were completely engrossed. It seemed to me to be designed by a committee of psychologists and educators. The link more or less confirms it. As for ``manipulation'', the creators admit that their intent is to portray certain things in a positive light. Whether right or wrong, the intent is there.

\hfill

\texttt{Daniel on 2010-10-24 at 22:58 said:}

Yes, of course, I see. So much of television is like this (explicitly intending to indoctrinate) that I thought perhaps there was something specific about that show. But the very ubiquity of the phenomenon makes your point exactly.

\hfill

\texttt{kadambari on 2010-10-31 at 11:48 said:}

``The final discovery is that of the unconscious in human motivation.''

This was recognized by ancient civilizations, that the unconscious is a potent force; however the ideal was always to overpower it by something higher and tame it, this marked the superior man, you see this in the Doric harmony in Greece and in Hindu mythology where the ideal is self-control and self-mastery, the overcoming of desires and the irrational. This is different from ``repression'' of these forces which is not self-mastery or self-growth. The self-mastery results in our being minimally influenced by them as they are placed under something higher; it has to come about as self-evolution and self-growth, cannot be a product of obedience to external commandments, and cannot come about by force. Repression comes about where the spirit does not dominate.

I think Freud was the biggest disaster as far as the treatment of the unconscious goes. When I was younger, I remember upon reading some of his works for the first time thinking deep down that there was something completely wrong with it, even perverse, but was too naive to exactly articulate what was wrong. Now I realize perhaps he is just a product of his own culture as exemplified in the Old Testament, where a God is to be obeyed and feared, and people brought under this culture exhibit the kinds of neurosis he deals with eventually as a cultural outcome? Perhaps his explanation is valid for only that cultural archetype and for those brought up under this archetype? However, this cultural archetype is gaining dominance as exemplified in the aesthetics of the movies, music and dance in popular culture, even our conceptions of beauty in the so called fashion magazines which are lacking any connection to higher aesthetic ideals.

Now there are different cultures and their treatment of the unconscious was quite different, but the Freudian explanation has come to dominate with disastrous consequences.

Hence today, in literature which generally reflects the culture of the times, importance is given to the irrational, the morbid, the perverse, which in higher civilizations is always placed under higher forces, and not given the sole importance and kept in their proper place.

Hence, one cannot feel that there is a certain ``sickness'' in popular dance, movies, music and such as most of it fails to appeal to anything higher in the human psyche. It is ultimately what a culture gives importance to; if the lowest drives and the irrational is thought worthy of veneration, then the cultural by-product as seen in the movies, dance and music and literature is just a reflection of that.

\end{sffamily}
\end{footnotesize}

